\section{Aplicación de la metodología y alcance de las conclusiones}
\label{bre:closing}

\subsection{Condiciones que deben acompañar a cada resultado}

El procedimiento comienza con una ecuación bien definida. Para cada caso
se especifican el campo, sus parámetros, las coordenadas y escalas, el
dominio y las superficies donde cambia su regularidad. En orden
fraccionario se añaden el operador, el orden de cada componente, el
terminal inferior y la inicialización. El argumento de existencia local
se aplica en el dominio utilizado; la existencia global exige una cota
adicional o se limita la conclusión al intervalo anterior a la salida.
Un cambio de unidades temporales modifica el segundo miembro mediante
la potencia $q$ de la escala, como se deduce del cambio de variable en
la integral de Caputo.

La construcción algebraica de una semilla requiere sustituirla en las
ecuaciones originales. En las eliminaciones se examinan por separado
los factores divididos y los denominadores nulos; se descartan los
equilibrios cuando se buscan PP. En un balance armónico se conservan el
signo de realimentación, el término medio y la fase de reconstrucción.
El residuo se calcula en las coordenadas físicas, incluyendo el efecto
de los armónicos descartados. Una raíz aproximada de las ecuaciones
proyectadas no establece la existencia de una solución de la ecuación
funcional completa.

La continuación debe llegar al campo objetivo. En una EDO, el estado final
de una etapa basta para comenzar la siguiente. En Caputo se conserva la
contribución de todas las etapas anteriores y se tratan las discontinuidades
temporales de la familia. Si después se utiliza el estado final como dato
de un PVI autónomo con memoria nula, se declara ese reinicio y se integra
su nueva trayectoria. El balance estacionario, la trayectoria durante la
continuación y la trayectoria autónoma de reinicio corresponden a tres
problemas relacionados, con condiciones iniciales distintas.

La validación numérica compara aproximaciones del mismo problema.
Refinar paso, ampliar horizonte y variar una perturbación inicial
responden a preguntas distintas. La comparación entre códigos tampoco
sustituye una cota de error si ambos comparten una inicialización
inconsistente. El residuo integral y los ejemplos exactos de control
permiten examinar el operador efectivamente aproximado. Para soluciones
sensibles, la separación puntual de trayectorias a tiempos largos se
complementa con retornos y observables definidos sobre ventanas
equivalentes; no se exige una coincidencia punto por punto incompatible
con la sensibilidad dinámica.

La identificación de un destino atractivo exige algo más que acotación
observada. Una región positivamente invariante puede contener varios
destinos. Un exponente variacional estimado positivo, una señal de banda
ancha o un estadístico cercano al régimen caótico aportan diagnósticos
complementarios, sin constituir por separado una demostración de caos.
La ausencia de contactos en un conjunto finito de condiciones iniciales
tampoco implica la exclusión de una vecindad abierta completa.

\subsection{Cobertura de las cuatro familias}

La Tabla~\ref{bre:closing:coverage} organiza el alcance de los cálculos.
La distinción entre tipo Chua y no Chua corresponde a la estructura del
campo; la distinción entre orden entero y fraccionario corresponde al
operador temporal. Una representación de Lur'e en un sistema no Chua no
cambia esa clasificación.

\begin{table*}[t]
\centering
\caption{Resultados disponibles y operaciones adicionales para una conclusión más fuerte.}
\label{bre:closing:coverage}
\small
\begin{tabularx}{\textwidth}{@{}L{0.14\textwidth}L{0.20\textwidth}Y Y@{}}
\toprule
Familia & Sistemas & Resultado establecido o calculado & Desarrollo adicional\\
\midrule
Entero tipo Chua & Chua PWL & Equilibrios, estabilidad, balance y semilla;
trayectorias, comparación entre implementaciones e indicadores finitos. &
Control de atracción y separación de vecindades completas para una prueba de ocultedad.\\
Entero no Chua & MAVPD, Kalman--Fitts, PLL &
Balance o demostración de su obstrucción; rutas alternativas, retornos e
indicadores. Resolución de PP y ausencia de PP según el sistema. &
Validación de los operadores de retorno y de sus dominios; regiones
invariantes que separen cuencas. En MAVPD, prueba de caos adicional a los exponentes estimados.\\
Fraccionario tipo Chua & PWL y arctangente &
Formulación causal, estabilidad sectorial, semillas y trayectorias con
memoria. Sondeos de reinicio; controles de paso donde se indican. &
Completar el refinamiento del caso PWL; controlar el espacio de perfiles,
la atracción y la exclusión de vecindades para fortalecer la clasificación.\\
Fraccionario no Chua & Lorenz generalizado, RF, Muñoz--Pacheco &
Equilibrios y PP/FPP--A; acotación y absorción de Lorenz. Integraciones
piloto de semillas Muñoz y diagnósticos específicos RF. &
Completar localización autónoma con memoria y controles por caso;
establecer atracción y cuencas. En RF, no trasladar resultados entre
valores distintos de $b$.\\
\bottomrule
\end{tabularx}
\end{table*}

Para Lorenz generalizado, la cota absorbente permite continuar todas las
soluciones de reinicio y restringe la región de búsqueda. No identifica
el conjunto límite al que llega cada semilla. El siguiente cálculo
consiste en integrar las semillas algebraicas bajo un mismo PVI,
comparar los destinos al refinar la discretización y analizar su
atracción. Para RF, cada conjunto de parámetros requiere sus propios
equilibrios, márgenes sectoriales, condiciones iniciales y controles de
existencia durante el horizonte usado. Para Muñoz--Pacheco, la ausencia
de equilibrios elimina una condición de exclusión local, pero no
demuestra acotación ni atracción: la trayectoria exacta no acotada sobre
$x=z=0$ impide inferir acotación global a partir de los pilotos.

El análisis topológico puede reducir esta parte del trabajo cuando se
construye un objeto al que aplicar un teorema. Un grado no nulo permite
probar un cero del mapa correspondiente; un retorno contractivo en un
dominio invariante permite probar un punto fijo atractivo; una
separatriz demostrada puede excluir el acceso desde una región. El
objeto, su dominio, la regularidad y la condición de frontera deben
verificarse antes de aplicar el resultado. En orden fraccionario, la
proyección puntual de una trayectoria no proporciona automáticamente
un mapa de retorno: es necesario conservar la memoria o justificar una
reducción funcional compatible.

Los desarrollos anteriores permiten ejecutar la búsqueda y reconocer el
alcance preciso de cada resultado. Las cuestiones abiertas se expresan
como operaciones matemáticas o numéricas concretas: construir un
dominio invariante, validar un retorno, controlar el defecto integral,
identificar una familia admisible de historias o excluir vecindades
completas. De este modo, la metodología distingue los pasos resueltos
de las hipótesis que todavía requieren demostración o cálculo.
